\section{Agent 定义与分类：感知 + 行动}

导读说明四份材料都指向 Agent，本章先把概念底座搭稳，回答一个具体问题：什么系统才算 Agent，而什么只是聊天机器人。视频描述给出的简洁定义是：Agent 是能够与环境进行交互的智能系统，具备感知能力和行动能力。感知包括观察环境、读取反馈、解析上下文；行动包括调用工具、生成输出、控制界面、修改变量。最小循环是“观察 → 决策 → 行动”。

\begin{figure}[H]
\centering
\includegraphics[width=0.96\textwidth]{agent-definition-loop.png}
\caption{Agent 最小定义：感知环境、决策、行动，再读取反馈。自制概念图，依据 00:02:00--00:14:50 对谈内容整理。}
\end{figure}

\begin{knowledgebox}{读图：Agent 和 Chatbot 的差别}
Chatbot 主要响应输入；Agent 要在环境中改变状态。只要系统开始调用工具、读反馈、调整计划，就从对话系统走向智能体系统。
\end{knowledgebox}

\subsection{四类 Agent}

有了最小定义之后，本节按环境和产品形态分类。节目把 Agent 分成 Coding Agent、Search Agent、Tool-use Agent 和 Computer-use Agent。这个分类不是严格学术分类，而是产品和环境分类：代码、搜索、工具、电脑界面，对应不同 affordance 和验证方式。

\begin{figure}[H]
\centering
\includegraphics[width=0.96\textwidth]{agent-types-map.png}
\caption{Agent 类型地图：Coding、Search、Tool-use、Computer-use 对应不同环境。自制概念图，依据 00:02:00--00:14:50 对谈内容整理。}
\end{figure}

\begin{knowledgebox}{术语消化：Agent 类型}
\begin{tabularx}{\textwidth}{p{0.22\textwidth}p{0.34\textwidth}X}
\toprule
类型 & 代表任务 & 难点 \\
\midrule
Coding Agent & 代码补全、重构、调试、PR 修改 & 代码库上下文、测试、长程修改。 \\
Search Agent & 检索、汇总、调研报告 & 信息源可信、引用、去重和 synthesis。 \\
Tool-use Agent & 调用 API、函数、数据库和业务工具 & 工具选择、参数生成、错误恢复。 \\
Computer-use Agent & 操作浏览器、GUI、跨应用流程 & 视觉 grounding、状态追踪、权限安全。 \\
\bottomrule
\end{tabularx}
\end{knowledgebox}

\subsection{本章小结}

Agent 的定义可以很简单，但落地非常复杂。不同 Agent 类型的本质差别，在于它们面对的环境、工具和反馈不同。

